\documentclass[11pt]{article}
\usepackage[utf8]{inputenc}
\usepackage[T1]{fontenc}
\usepackage[margin=0.9in]{geometry}
\usepackage{amsmath,amssymb,amsfonts,mathtools,bm, longtable, setspace}
\usepackage{booktabs,array,enumitem,xcolor,colortbl, makecell}
\usepackage[hidelinks,hypertexnames=false]{hyperref}
\usepackage{caption}
\captionsetup[longtable]{
  justification=centering,
  singlelinecheck=true
}
\setlength{\LTcapwidth}{\textwidth}
\setlist{nosep,leftmargin=1.4em}
\newcommand{\clip}{\operatorname{clip}_{[0,1]}}
\newcommand{\Real}{\mathbb{R}}
\newcommand{\half}{\tfrac{1}{2}}
\newcommand{\Nn}{\mathcal{N}}
\DeclareMathOperator*{\argmin}{arg\,min}

\title{\textbf{ODD Protocol for State-Dependent Participation Diffusion on Adaptive Influence Networks}}
\author{Hudson Dong\thanks{Class of 2027, Stuyvesant High School, 345 Chambers St, New York, NY 10282. Email: \texttt{hdong70@stuy.edu}}}
\date{}
\begin{document}
\setstretch{1.1}
\maketitle
\begin{abstract}
This document describes the agent-based model in \textit{State-Dependent Participation Diffusion on Adaptive Influence Networks} using the ODD protocol. The model represents recurrent collective participation on a weighted, directed, adaptive influence network. Agents may participate, broadcast, become inactive, and later return. Participation and social interaction change motivation, commitment, fatigue, and influence relationships, while external events, capacity, and temporary scaffolding alter the conditions under which agents participate, leave, or return. The implementation supports stochastic, deterministic expected-value, and ensemble simulation. The specification records the update sequence, initialization rules, parameter values, and sub-models used in the accompanying implementation. 
\end{abstract}

\section{Overview}
\subsection{Purpose and Patterns}
The model studies recurrent collective participation on an adaptive influence network. Unlike one-time adoption models, it follows repeated participation, temporary inactivity, and return as internal states and influence relationships change. Voluntary student activities provide the motivating example.

The central question is whether temporary activation develops into sustained participation or disappears after external stimulation is withdrawn. The main patterns of interest are collapse after a temporary spike and continued high participation after external stimulation ends. The model also tracks threshold sensitivity and path dependence generated by accumulated commitment. Turnover and reactivation provide secondary population-level outcomes, while repeated co-participation and broadcasting change the influence network.

The reference experiments compare collapse-prone, persistent, and intervention-supported regimes. The control analysis varies event pulses, visibility, capacity, and temporary scaffolding.

\subsection{Entities, State Variables, and Scales}
The model contains $n$ agents indexed by $i\in V$ and a weighted, directed influence network:
\[
G_t=(V,E,W_t).
\]
A directed edge $j \xrightarrow{} i$ represents a channel through which agent $j$ may influence agent $i$. The effective weight of this channel is
\[
\omega_{ij}(t)=\clip[v_{ij}(t)s_{ij}(t)],
\]
where $v_{ij}\in [0,1]$ denotes visibility and $s_{ij}(t)\in[0,1]$ denotes persuasiveness. The node set remains fixed during each simulation run. However, activity status and influence weights may change over time.

Each agent carries three core latent states: motivation $m_i(t)\in[0,1]$, commitment $L_i(t)\in [0,1]$, and fatigue $F_i(t)\in[0,1]$. Optional modulators include affect $E_i(t)\in[0,1]$, saturation $S_i(t)\in[0,1]$, and stability gain $\Phi_i(t)\in[0,1]$. Motivation is the fast state and governs immediate behavioral propensity. Commitment and fatigue change more slowly and affect later motivational decay. When enabled, the modulators adjust action propensity and the response to reinforcement, social load, and recent variation. 

The activity variable $act_i(t)$ distinguishes active agents from those who are temporarily inactive. In stochastic simulations,
\[
act_i(t)\in\{0,1\},
\]
whereas in deterministic expected-value simulations $act_i(t)\in[0,1]$ is propagated as an expected active fraction. The two observable behavioral variables are participation $a_i(t)$ and broadcasting $p_i(t)$. In the motivating implementation, they correspond to attendance and posting. More generally, they may represent other costly and low-cost forms of recurrent action. Capacity adjustment converts intended participation $a_i(t)$ into realized participation $\tilde{a}_i(t)$. Aggregate realized participation is
\[
A(t)=\sum_{i=1}^n\tilde{a}_i(t).
\]
Agent-specific parameters govern action thresholds, decay, peer response, reinforcement, fatigue, dropout, and reactivation. The model also includes several externally supplied covariates. The availability factor $c_i(t)\in[0,1]$ determines how much of an event input reaches agent $i$. The quantity $T_i(t)\ge0$ represents the available engagement opportunity, while $\bar{y}_i>0$ caps credited engagement. The role-salience variable $s_i^{\text{role}}(t)$ scales the commitment generated through socially reinforced participation. Finally, $r^+(t)\in[0,1]$ represents favorable collective feedback. These quantities are scenario inputs or observable covariates, not latent states.

The population-level quantities are the viability threshold $N_{min}$, the run-level capacity $K$, the event schedule $\{(\tau_k,A_k)\}$, the run-level visibility multiplier $\gamma_v$, and the time-dependent scaffolding multiplier $\gamma_B(t)$. In the current implementation, $K$ and $\gamma_v$ remain constant within a simulation run. Event input and scaffolding, however, may vary across periods.

Time is discrete, with model periods
\[
t=0, 1, \ldots T-1.
\]
Displayed Week $w$ corresponds to model period $t=w-1$. Because actions are generated before the event input updates motivation, an event scheduled for Week $w$ first affects participation in Week $w+1$. For the 18-period experiments, the terminal four-week window is $t=14,\dots, 17$. One model period represents one week in the accompanying implementation. Core latent states and normalized influence variables are projected onto $[0,1]$ after each update.

\subsection{Process Overview and Scheduling}
Updates are synchronous: all period-\(t\) quantities are evaluated before any period-\(t+1\) state is written. Within-period outcomes depend only on the states and network available at the start of the period. Each period proceeds from current states and network weights to actions, realized participation, state transitions, and network updating.

One period proceeds as follows:
\begin{enumerate}
    \item read the period-\(t\) agent states, network weights, exogenous inputs, and controls;
    \item calculate the motivation entering the action rule and the participation and broadcasting probabilities;
    \item draw stochastic actions or substitute their deterministic expected
values;
    \item apply the capacity adjustment to obtain \(\tilde a_i(t)\) and \(A(t)\);
    \item calculate peer-motivation input and weighted neighbor participation using the period-\(t\) network;
    \item calculate external input, internal reinforcement, effective decay,   dropout hazard, and reactivation probability;
    \item update motivation, commitment, fatigue, optional modulators, and activity status synchronously;
    \item update \(v_{ij}(t+1)\) from period-\(t\) participation, broadcasting, and the persistent channel term;
    \item calculate \(cred_j(t+1)\), \(s_{ij}(t+1)\), and \(\omega_{ij}(t+1)\) for the next period; 
    \item record period-\(t\) actions and aggregate outputs.
\end{enumerate}

The updated visibility, credibility, persuasiveness, and effective influence enter only the next period's action rule. The snapshot labeled \(t\) contains the period-\(t\) states and influence weights together with the actions and aggregate participation generated from them. Values computed for \(t+1\) first appear in snapshot \(t+1\).

Stochastic, deterministic expected-value, and ensemble runs use this same transition order; only the treatment of action draws and randomized realizations differs.

\section{Design Concepts}
\subsection{Basic principles}
The model distinguishes one-time activation from recurrent participation. Participation is generated by a fast motivational state and modified by slower commitment and fatigue states. Realized participation feeds back into these states, changing both the probability and persistence of future behavior.

The influence structure is also endogenous. Visibility and persuasiveness change through public signaling, co-participation, commitment, and recent participation. Events, capacity, visibility, and scaffolding modify specified terms in the state transitions.

\subsection{Emergence}
Aggregate participation and its persistence or collapse emerge from individual state transitions and directed interaction. Turnover and intervention responses arise from the same process. Agents are not assigned fixed types such as “persistent” or “inactive.” Instead, these patterns arise from differences in thresholds, decay rates, reinforcement parameters, network positions, and prior state histories.

The main emergent contrast is between a temporary peak followed by collapse and a high-activity trajectory sustained after the initial stimulus. The piecewise reinforcement function creates sensitivity to the viability threshold. Path dependence and hysteresis can arise because commitment accumulated during earlier participation lowers later motivational decay

\subsection{Adaptation}
Agents adapt through changes in their internal states and influence relationships. Motivation responds quickly to peer input, external events, and realized reinforcement. Participation and collective feedback build commitment, while engagement and social exposure build fatigue; inactivity permits partial recovery.

The network adapts through visibility and credibility. Co-participation and broadcasting change exposure, while commitment and recent participation change persuasiveness.

\subsection{Objectives}
Individual agents have no explicit utility function or forward-looking objective. Their actions follow stochastic or deterministic threshold responses to current states.

The reverse planner is external to the agents. It searches over events, visibility, capacity, and temporary scaffolding to meet a terminal-window participation target at minimum control cost. The planner varies controls while leaving the transition equations unchanged.

\subsection{Learning and Prediction}
Agents do not learn strategically or estimate a model of their environment. Past experience affects behavior through accumulated state variables. Previous participation affects commitment and credibility, while repeated interaction changes visibility. These state variables provide memory without explicit belief updating or statistical learning.

Forward prediction occurs in the simulator and reverse planner, not within the agents. The planner evaluates each candidate control by propagating the model over the remaining horizon. The resulting projections are conditional on the selected parameters and initial state.

\subsection{Sensing}
Agent $i$'s action depends on its internal states, activity status, and exogenous inputs. Social information enters through weighted incoming influence: peer motivation is transmitted through  $\omega_{ij}(t)$, and neighbor participation through $\bar{a}_{N(i)}(t)$.

Neighbors’ motivation is not observed directly; it enters through the weighted peer signal in Equation (2). Event exposure depends on availability, event visibility, affect, and saturation. 

\subsection{Interaction}
Interactions are directed and weighted. The influence of agent $j$ on agent $i$ depends on the product of visibility $v_{ij}(t)$ and persuasiveness $s_{ij}(t)$. Co-participation changes mutual visibility, while broadcasting changes the sender’s directed visibility.

The network transmits motivational and behavioral signals, and resulting actions may later alter the network. Inactive agents do not participate or broadcast, but external input or active neighbors may reactivate them.

\subsection{Stochasticity}
Randomness enters through network generation, parameter and state initialization, action draws, shocks, dropout, and reactivation. In stochastic runs, participation and broadcasting are Bernoulli variables conditional on current action probabilities and activity status.

Deterministic expected-value runs replace these draws with conditional expectations and propagate expected activity status. The deterministic mode is a plug-in approximation: expected actions and activity fractions are inserted directly into the nonlinear transition equations. The reported deterministic ensembles still resample the network, heterogeneous parameters, and initial states.

\subsection{Collectives}
The network generator assigns agents to initial communities that create clustered ties and cross-community bridges. These communities shape exposure but have no separate states or objectives. Population-level conditions are represented by aggregate participation $A(t)$ and the viability threshold $N_{\min}$.

\subsection{Observation}
The primary observable outputs are individual participation, broadcasting, realized participation, activity status, and aggregate participation. The simulator also records motivation, commitment, fatigue, the optional modulator states, and effective influence weights. Motivational reach is the fraction of agents satisfying $m_i(t)>0.4$.

Deterministic and ensemble analyses report means, terminal-window averages, and empirical quantiles across randomized realizations. Public-body dots are display elements outside the modeled node set $V$.

\section{Details}
\subsection{Initialization}
\subsubsection{Network Initialization}
Each simulation begins with \(n\) agents and a directed, clustered influence network. Let \(g\) denote the target friend-group size. The number of communities is
\[
C=\max\left\{2,\operatorname{round}\left(\frac{n}{g}\right)\right\}.
\]
With \(n=36\) and \(g=8\), the reference network contains five communities of sizes \(8,7,7,7,\) and \(7\). Agents are assigned by index,
\[
\mathrm{comm}_i=i\bmod C.
\]
Community labels are fixed by index, while angular positions and network ties vary with the random seed.

Agent \(i\) receives angular position
\[
\theta_i
=
\frac{2\pi\,\mathrm{comm}_i}{C}+U_i,
\qquad
U_i\sim
\operatorname{Uniform}\left(
-\frac{0.8\pi}{C},
\frac{0.8\pi}{C}
\right),
\]
and dyadic similarity is
\[
\mathrm{sim}_{ij}
=
\exp\left(
-\frac{d_{ij}^{2}}{0.6}
\right).
\]
For each receiver \(i\), the number of within-community incoming ties is
\[
k_i^{\mathrm{in}}
=
\operatorname{round}(d_{\mathrm{in}})+U_i^{\mathrm{in}},
\qquad
U_i^{\mathrm{in}}
\sim
\operatorname{DiscreteUniform}\{0,1,2\}.
\]
With \(d_{\mathrm{in}}=4\), each receiver forms incoming ties from the four to
six most similar agents in its community. The number of cross-community bridge ties is
\[
k_i^{\mathrm{out}}
=
\left\lfloor d_{\mathrm{out}}\right\rfloor+Z_i,
\qquad
Z_i\sim
\operatorname{Bernoulli}\left(
d_{\mathrm{out}}-\left\lfloor d_{\mathrm{out}}\right\rfloor
\right).
\]
With \(d_{\mathrm{out}}=1.5\), each receiver obtains one or two uniformly
sampled cross-community senders with equal probability. Self-loops and duplicate directed edges are prohibited. Reciprocal ties are allowed and generated independently.

For each standard within-community or bridge edge,
\[
v_{ij}^{\mathrm{mix}}(0)
=
\operatorname{clip}_{[0,1]}
\left[
b_{ij}
+
0.30\,\mathrm{sim}_{ij}
+
U_{ij}^{v}
\right],
\]
with \(b_{ij}\) and \(U_{ij}^{v}\) defined as above.

The implementation sets
\[
\mathrm{online}_{ij}(0)
=
\mathrm{offline}_{ij}(0)
=
v_{ij}^{\mathrm{mix}}(0),
\qquad
\pi_o=\pi_f=0.5,
\]
and holds both channel components fixed. Hence,
\[
v_{ij}^{\mathrm{mix}}(t)
=
v_{ij}^{\mathrm{mix}}(0)
\qquad
\text{for all }t.
\]
Initial adaptive visibility is
\[
v_{ij}(0)
=
\operatorname{clip}_{[0,1]}
\left[
0.5v_{ij}^{\mathrm{mix}}(0)+0.1
\right].
\]
Because no pre-simulation participation history is supplied,
\[
\bar{a}_{j}^{\mathrm{pre}}=0,
\qquad
\mathrm{cred}_{j}(0)=0.5L_j(0).
\]
Initial persuasiveness and effective influence are
\[
s_{ij}(0)
=
\operatorname{clip}_{[0,1]}
\left[
0.30
+
0.40\,\mathrm{sim}_{ij}
+
0.30\,\mathrm{cred}_{j}(0)
\right],
\qquad
\omega_{ij}(0)
=
\operatorname{clip}_{[0,1]}
\left[
v_{ij}(0)s_{ij}(0)
\right].
\]
The master seed is \(q_0=11\), and realization $r=0,\ldots,R-1$ uses
\[
q_r=11+9973r.
\]
Each realization resamples angular positions, directed ties, heterogeneous parameters, and initial states; scenario controls remain fixed. 

\subsubsection{Agent States and Parameters}
I draw the initial motivation \(m_i(0)\), commitment \(L_i(0)\), and fatigue
\(F_i(0)\) from the regime-specific distributions in Table B2a and clip each
draw to \([0,1]\). When the optional modulators are enabled, I initialize
\[
E_i(0)=\bar E_i,
\qquad
S_i(0)=0,
\qquad
\Phi_i(0)=1,
\]
together with
\[
\bar m_i(0)=m_i(0),
\qquad
\operatorname{var}_i(0)=0.
\]
All agents begin active:
\[
act_i(0)=1.
\]
Agent-specific parameters govern action thresholds, response slopes,
motivational decay, peer influence, reinforcement, fatigue, dropout,
reactivation, visibility, and persuasiveness. Unless a quantity is defined as a dynamic state or time-dependent input, its sampled value remains fixed throughout the realization.

The reference experiments use two parameter regimes. The collapse-prone regime combines faster motivational decay, weaker reinforcement and peer transmission, and a higher viability threshold. The favorable regime uses slower decay, stronger reinforcement and peer transmission, and a lower
viability threshold. These regimes are parameter configurations rather than fixed agent types.

Tables B1 to B3 report the parameter values, initialization distributions, and sampling status used in the implementation. A quantity marked "Fixed'" remains constant across the stated regime. A quantity marked "Randomized once" is sampled during initialization and then held fixed within that realization. Quantities marked "Dynamic" are updated by the model equations or supplied as time-dependent inputs.

\subsubsection{Controls and Reference Scenarios}
I fix the external controls before each forward simulation unless I use a receding-horizon procedure. An event schedule specifies the event times $\tau_k$ and amplitudes $A_k$. In the current implementation, I hold the visibility multiplier $\gamma_v$ and capacity $K$ constant within each simulation run. Temporary scaffolding may vary across periods through $\gamma_B(t)$. Outside the scaffolding window, $\gamma_B(t)=1$. The parameter $s_B$ determines the additive below-quorum support supplied while scaffolding is active.

The computational experiments reported in the associated paper use $n=36$ agents over an 18-period horizon. The model periods are $t=0, \ldots, 17$, which are displayed as Weeks 1--18. Displayed Week $w$ corresponds to $t=w-1$. Event and scaffolding labels in Table 1 follow the displayed-week convention. Three reference scenarios are compared:

\begin{table}[h]
    \centering
    \begin{tabular}{|c|c|c|c|c|c|}
    \hline
         Scenario&  Regime&  Events&  Visibility&  Capacity& Scaffolding\\
    \hline
         A&  Collapse-prone&  \makecell{Week 3\\amplitude 0.60}&  $\gamma_v=1.0$&  $K=14$& None\\
    \hline
         B&  Favorable& \makecell{Weeks 3, 7, 11\\amplitude 0.35 each}&  $\gamma_v=1.4$&  $K=18$& None\\
    \hline
         C&  Collapse-prone&  \makecell{Weeks 2, 5, 8, 12\\amplitude 0.45 each}&  $\gamma_v=3.0$& $K=20$& \makecell{Weeks 2--9\\$\gamma_B=1.4$}\\
    \hline
    \end{tabular}
    \caption{Reference control configurations used in the computational experiments}
\end{table}
Scenario C is the rescue configuration used in the 18-period experiment.

In the Figure 1 parameter sweep, I hold the Scenario A event schedule, capacity, and scaffolding settings fixed. I vary only $N_{min}$ and $\gamma_v$.

Figures 1 and 2 use deterministic expected-value actions, with $a_i(t)$, $p_i(t)$, and activity status propagated through their conditional expectations.

In the reported reference experiments, I activate the affect, saturation, and stability modulators with the nonzero coupling values listed in Table B1. I set the idiosyncratic motivation shock to $\epsilon_i(t)=0$ in the deterministic experiments. Each realization resamples the network, heterogeneous parameters, and initial states while retaining the same scenario controls, horizon, population size, and regime. Each cell in the viability-visibility sweep averages 20 independent realizations. The trajectory comparison averages $R=250$ independent realizations for each reference scenario. 

The realization seeds follow the master-seed rule in Table B3, and the computational materials record the seed sequence and implementation version used for the figures.

\subsection{Input Data}
The current implementation uses no empirical input data. I generate the networks, initial states, and heterogeneous parameters synthetically from the distributions of the selected regime. The user or a reference scenario supplies the event schedule, visibility amplification, capacity, scaffolding, simulation horizon, random seed, and ensemble size.

In synthetic runs, I generate the externally supplied agent-level quantities $c_i(t)$, $T_i(t)$, $s_i^{role}(t)$, and the population-level signal $r^+(t)$ from scenario settings. An empirical implementation would require these inputs to be estimated or constructed from longitudinal participation and network data.

\subsection{Submodels}
I use the following definitions throughout the submodels:
\[
\sigma(x)=\frac{1}{1+e^{-x}}, \qquad \clip(x)=\min\{1, \max\{0, x\}\}, \\
\left[ x \right]_+ = \max\{x,0\}, \qquad \left[ x \right]_- = \max\{-x,0\},
\]
and
\[
\mathbf{1}\{\mathcal{C}\}=
\begin{cases}
    1, & \text{if condition  }  \mathcal{C} \text{ is true,} \\
    0, & \text{otherwise.} \\
\end{cases}
\]
The constant $\epsilon_0>0$ prevents division by zero. The term $\epsilon_i(t)$ is a zero-mean idiosyncratic motivation shock with the distribution reported in Table B2; In the deterministic experiments, however, I set $\epsilon_i(t)=0$.

\subsubsection{Actions and Realized Participation}
Emotion-shifted motivation is
\[
\widehat{m}_i(t)=\clip\bigg[m_i(t)+\beta_E\Big(E_i(t)-\frac{1}{2}\Big)\bigg]
\]
Participation and broadcasting probabilities are
\[
P_i^a(t)=\sigma(\kappa_i^a[\widehat{m}_i(t)-\theta_i^a]),
\]
\[
P_i^p(t)=\sigma(\kappa_i^p[\widehat{m}_i(t)-\theta_i^p]).
\tag{1}
\]
In stochastic simulations,
\[
Z_i^a(t) \sim \text{Bernoulli}(P_i^a(t)), \qquad Z_i^p(t) \sim \text{Bernoulli}(P_i^p(t)),
\]
and
\[
a_i(t)=act_i(t)Z_i^a(t), \qquad p_i(t)=act_i(t)Z_i^p(t).
\]
In deterministic expected-value simulations,
\[
a_i(t)=act_i(t)P_i^a(t), \qquad p_i(t)=act_i(t)P_i^p(t).
\]
I convert intended participation into realized participation through the capacity constraint:
\[
\widetilde{a}_i(t)=\frac{a_i(t)}{1+\frac{\sum_ja_j(t)}{K}}=a_i(t)\frac{K}{K+\sum_ja_j(t)}.
\tag{14}
\]
Aggregate realized participation is:
\[
A(t)=\sum_{i=1}^n \widetilde{a}_i(t).
\]

\subsubsection{Peer Input, Motivation, and Decay}
I define the weighted local participation aggregate as:
\[
\bar{a}_{N(i)}(t)=\frac{\sum_j\omega_{ij}(t)\widetilde{a}_j(t)}{\sum_j\omega_{ij}(t)+\epsilon_0}.
\]
The peer-motivation signal received by agent $i$ is
\[
X_i(t)=\sum_j\omega_{ij}(t)b_j\tanh(\kappa_gm_j(t)),
\]
\[
\widetilde{X}_i(t)=\frac{X_i(t)}{1+\rho_iX_i(t)}.
\tag{2}
\]
I update motivation according to
\[
m_i(t{+}1)=\clip\!\Big[(1-\alpha_i(t))m_i(t)+G(\Phi_i(t))\big(w_i^P(1-S_i(t))\tilde X_i(t)+c_i(t)R_i(t)+B_i(t)+\epsilon_i(t)\big)\Big]
\tag{3}
\]
where
\[
G(\Phi)=1+\kappa_\Phi(1-2\Phi), \qquad \kappa_\Phi\in[0,1).
\]
The effective fatigue entering the decay equation is:
\[
\widehat{F}_i(t)=\clip\Big[F_i(t)+\kappa_{EF}(\frac{1}{2}-E_i(t))\Big].
\]
The endogenous motivational decay rate is then:
\[
\alpha_i(t)=\clip\Big[\alpha_{0,i}(1-\eta_L[1-\kappa_{SL}S_i(t)]L_i(t))\big(1+\eta_F\widehat{F}_i(t)\big)\Big].
\tag{4}
\]

\subsubsection{Commitment, Fatigue, and Affective Modulation}

I use separate multipliers to represent the effects of affect on fatigue and positive reinforcement:
\[
\nu_F(E)=1+\beta_F(\frac{1}{2}-E),
\]
\[
\nu_R(E)=1-\beta_R(\frac{1}{2}-E),
\]
with
\[
0\le \beta_F, \beta_R \le 2.
\]
I update commitment as:
\begin{equation}\tag{5}\label{eq:loyalty}
\begin{aligned}
L_i(t{+}1)=\clip\!\Big[L_i(t)+\lambda_L\big(&\phi_1\tilde a_i(t)+\phi_2s_i^{\mathrm{role}}(t)(1-\kappa_{SL}S_i(t))\bar a_{N(i)}(t)\\
&+\phi_3\nu_R(E_i(t))r^+(t)-L_i(t)\big)\Big].
\end{aligned}
\end{equation}
I update fatigue as:
\begin{equation}\tag{6}\label{eq:fatigue}
\begin{aligned}
F_i(t{+}1)=\clip\!\Big[F_i(t)+\lambda_F\big(&\nu_F(E_i(t))\psi_1y_i(t)+\nu_F(E_i(t))\psi_2\bar a_{N(i)}(t)\\
&+\psi_3(1-\tilde a_i(t))-F_i(t)\big)\Big].
\end{aligned}
\end{equation}
\subsubsection{External and Internal Reinforcement}
I represent external stimulation as a sequence of exponentially decaying event pulses:
\[
R_i(t)=\kappa_R[1-S_i(t)]\nu_R(E_i(t))V_i^{ev}\sum_{k:\tau_k\le t}A_ke^{-\lambda_R(t-\tau_k)}, \qquad \kappa_R=3.4.
\tag{7}
\]
Internal reinforcement follows the piecewise function:
\begin{equation}\tag{8}\label{eq:internal}
B_i(t)=\nu_R(E_i(t))
\begin{cases}
\gamma_B(t)\beta_i\big(1-e^{-\delta_i y_i(t)}\big)-\chi_i\mathbf 1_{\{y_i(t)=0\}}, & A(t)\ge N_{\min},\\[3pt]
-\chi_i^{\mathrm{viab}}+[\gamma_B(t)-1]s_B, & A(t)<N_{\min}.
\end{cases}
\end{equation}
The corresponding reinforcement half-saturation point is
\[
y_{i,1/2}=\frac{\ln2}{\delta_i}.
\]
Above quorum, scaffolding amplifies productive reinforcement. Below quorum, it offsets part of the viability penalty without amplification of the nonparticipation penalty.

\subsubsection{Adaptive Influence Weights}
I define effective influence as:
\[
\omega_{ij}(t)=\clip(v_{ij}(t)s_{ij}(t)).
\tag{9}
\]
The persistent channel mixture is defined as:
\[
v_{ij}^{mix}(t)=\pi_o\text{online}_{ij}(t)+\pi_f\text{offline}_{ij}(t), \qquad \pi_o+\pi_f=1.
\tag{12}
\]
I calculate persuasiveness as:
\[
s_{ij}(t)= \clip\big[\zeta_i^0+\zeta_i^1\mathrm{sim}_{ij}+\zeta_i^2\mathrm{cred}_j(t)\big],\tag{10}\label{eq:persuade}
\]
where
\[
\mathrm{cred}_j(t)=w_1^cL_j(t)+w_2^cH_j(t),\tag{11}\label{eq:cred}
\]
\[
d_t=\min\{\Delta,t\},
\]
and
\[
H_j(t)=
\begin{cases}
0, & t=0, \\
\frac{1}{d_t}\sum_{\tau=t-d_t}^{t-1}\tilde a_j(\tau) & t\ge1.
\end{cases}
\]
Credibility at time $t$ uses $L_j(t)$ and at most the the preceding $\Delta$ realized-participation observations. After the period-$t$ states and participation history are updated, Equations (10)--(11) determine $s_{ij}(t+1)$ and $\omega_{ij}(t+1)$ for the next behavioral update.

The visibility multiplier acts through the sender’s effective expressive reach:
\[
V_{j,\text{eff}}^{ex}(t)=\clip\big[\gamma_vV_{j,0}^{ex}\big].
\]
I then update visibility as:
\begin{equation}\tag{13}\label{eq:vupdate}
v_{ij}(t{+}1)=\clip\!\Big[(1-\mu)v_{ij}(t)+\mu\big(\eta_1\tilde a_i(t)\tilde a_j(t)+\eta_2V^{ex}_{j,\text{eff}}p_j(t)\big)+\mu_{\mathrm{mix}}\big(v_{ij}^{\mathrm{mix}}(t)-v_{ij}(t)\big)\Big].
\end{equation}
The visibility control changes the opportunity for exposure, but persuasiveness remains unmodified.

\subsubsection{Attrition and Reactivation}
I define the dropout hazard as
\[
h_i(t)=h_{0,i}\exp\big[-\gamma_1m_i(t)-\gamma_2L_i(t)+\gamma_3F_i(t)\big].
\tag{15}
\]
Conditional on agent $i$ being active,
\[
d_i(t)\sim \text{Bernoulli}\big(1-e^{-h_i(t)}\big).
\]
For an inactive agent, the reactivation probability is:
\[
P_i^{react}(t)=1-\exp\big[-\xi_1R_i(t)-\xi_2\bar{a}_{N(i)}(t)\big].
\]
\[
r_i(t)\sim \text{Bernoulli}\big(P_i^{react}(t)\big).
\tag{16}
\]
In stochastic simulations, I update activity status as:
\[
act_i(t+1)=act_i(t)[1-d_i(t)]+[1-act_i(t)]r_i(t).
\]
In deterministic expected-value simulations, I use:
\[
act_i(t+1)=act_i(t)e^{-h_i(t)}+[1-act_i(t)]P_i^{react}(t).
\]

\subsubsection{Optional Modulators}
I update affect according to:
\begin{equation}\tag{17}\label{eq:emotion}
E_i(t{+}1)=\clip\!\Big[(1-\lambda_E)E_i(t)+\lambda_E\bar E_i+\zeta_E[B_i(t)]_+-\zeta'_E[B_i(t)]_-+\kappa_E\xi_i^E(t)\Big].
\end{equation}
I define the incoming social load as:
\begin{equation}\tag{18a}\label{eq:saturation}
\ell_i(t)=\sum_j\omega_{ij}(t),
\end{equation}
I then update saturation as:
\begin{equation}\tag{18b}
S_i(t{+}1)=act_i(t)\Big[(1-\lambda_S)S_i(t)+\lambda_S\frac{\ell_i(t)}{\ell_i(t)+s_0}\Big].
\end{equation}
I calculate the moving mean and variance of motivation as:
\begin{equation}\tag{19a}
\begin{aligned}
\bar m_i(t)&=(1-\lambda_v)\bar m_i(t{-}1)+\lambda_vm_i(t),
\end{aligned}
\end{equation}

\begin{equation}\tag{19b}
\begin{aligned}
\mathrm{var}_i(t)&=(1-\lambda_v)\mathrm{var}_i(t{-}1)+\lambda_v(m_i(t)-\bar m_i(t))^2,\\
\end{aligned}
\end{equation}
The resulting stability gain is
\[
\Phi_i(t)=\frac{1}{1+\kappa_{\text{var}}\mathrm{var}_i(t)}.
\tag{19c}
\]
\subsubsection{Reverse Planning and Ensemble Evaluation}
I use $u$ to represent event timing and amplitudes, the run-level visibility multiplier $\gamma_v$, the run-level capacity $K$, and the time-dependent scaffolding schedule $\gamma_B(t)$. In the current implementation, $\gamma_v$ and $K$ are constant within each candidate simulation run. The planner searches for a low-cost control that satisfies the terminal-window target:
\[
\overline A_{T-w:T-1}(u)=\frac1w\sum_{t=T-w}^{T-1}A(t)\ge A^*.
\tag{20}
\]
A representative additive cost is
\[
C(u)=\kappa_{E,0}N_E+\kappa_{E,1}\sum_{k=1}^{N_E} A_k+\kappa_v(\gamma_v-1)+\kappa_B\sum_{t}[\gamma_B(t)-1]+\kappa_K[K-K_0]_+.
\]
Here, $N_E$ is the number of scheduled events, and $K_0$ is the uncontrolled baseline capacity. Table B1 reports the nonnegative control-cost coefficients $\kappa_{E,0}$, $\kappa_{E,1}$, $\kappa_B$, and $\kappa_K$. Because $\gamma_v$ and $K$ are run-level controls, their costs enter only once per simulation run. Scaffolding, however, is time dependent. Its cost is therefore summed across the periods in which it is active.

I use a bounded multi-start heuristic. Let $m$ denote the maximum permitted number of events. The heuristic begins by constructing three event-timing templates. The first template places maximum-amplitude events at displayed weeks
\[
w_k=1+\text{round}\bigg(\frac{k\min\{T-2,12\}}{\max\{1,m\}}\bigg), \qquad
k=0,\ldots,m-1.
\]
The second template distributes the same number of maximum-amplitude events evenly from the displayed-week $0.10(T-1)$ to $0.70(T-1)$. The third uses the interval from $0.25(T-1)$ to $0.80(T-1)$. Event weeks are rounded and clipped to the admissible range $1,\ldots,T-1$. 

I combine each event-timing template with a bounded set of capacity candidates. Let $K_{\max}$ denote the user-specified capacity ceiling. The capacity grid begins at $\min\{K_0, K_{\max}\}$ and increases in increments of two. I also include $K_{\max}$ when it does not fall exactly on this grid. Under the default values $K_0=14$ and $K_{\max}=20$, the tested capacities are
\[
K\in\{14,16,18,20\}.
\]
All starting configurations use the permitted visibility ceiling. When scaffolding is allowed, I set $\gamma_B=2.0$ over displayed Weeks 1-6 and truncate this window when the simulation horizon is shorter. When scaffolding is not allowed, I set $\gamma_B(t)=1$ throughout the candidate run.

I first evaluate the resulting starting configurations on the deterministic reference realization. If none satisfies the terminal target, the procedure reports that no feasible plan was found. It also displays the starting configuration with the highest terminal-window participation.

I order feasible starting configurations by control cost, from lowest to highest. If two configurations have the same cost, I evaluate the one with higher terminal-window participation first. Subject to the remaining global evaluation budget, I apply greedy control reduction to the feasible starts in this order. The global budget may be exhausted before every feasible start receives a complete descent. The procedure retains the lowest-cost feasible result found among the search paths that it actually evaluates.

When no immediate cost-reducing move remains feasible, the search tests one-period shifts in the existing event times. If enough of the evaluation budget remains, it may also evaluate up to five seeded perturbation starts. A perturbation may strengthen existing controls or add an event at another admissible week before greedy descent is resumed. Scenario C in Table 1 is a hand-specified rescue benchmark rather than an output of the heuristic planner.

Candidate controls are evaluated on the deterministic reference realization with seed 11. The planner enforces the terminal-window target without a downward feasibility tolerance. One global budget of at most 400 candidate-control evaluations covers the starting configurations, feasibility checks during greedy descent, event-timing changes, and perturbation searches. Each individual descent is limited to at most 60 passes. Therefore, some feasible starts may not receive a complete descent, and the procedure may evaluate fewer than five perturbation starts. The perturbation generator uses seed 12. Event amplitudes are reduced in increments of $0.05$, visibility in increments of $0.25$, scaffolding strength in increments of $0.15$, and scaffolding duration in one-period increments. Capacity candidates are separated by increments of two agents, with the user-specified ceiling included. When $\kappa_K>0$, the descent procedure may also reduce capacity by two agents. It accepts this change only if the stated cost decreases and feasibility is preserved. Under the reported setting $\kappa_K=0$, the control cost is unaffected by capacity, but the different capacity values remain distinct feasibility candidates.

The default reverse-planning interface uses $n=36$, $T=18$, the collapse-prone regime, terminal target $A^*=9$, a maximum of four events, an event-amplitude ceiling of $0.60$, a visibility ceiling of $\gamma_v=4.0$, a capacity ceiling of $K=20$, and permission to use temporary scaffolding. The default interface searches within these bounds; changing them may alter both feasibility and the plan returned by the heuristic.

I run the planner on one fixed deterministic reference realization. I then evaluate robustness separately over randomized ensembles. For scalar output $Z(t)$, the ensemble mean is
\[
\bar{Z}_R(t)=\frac{1}{R}\sum_{r=1}^RZ^{(r)}(t).
\]
Empirical quantiles summarize the variation across network, initial-state, and parameter realizations. A failed search indicates only that no feasible schedule was found within the evaluated candidates.

\appendix
\section{Visualization and Interface Specification}
\subsection{Motivation Color Map}
I use node fill color to represent motivation $m_i(t)$. The display uses a fixed plasma color map defined by six anchor colors placed at levels $\ell_k=k/5$.

\begin{table}[h]
    \centering
    \begin{tabular}{|l|l|l|l|l|}
    \hline
        $k$& Level&  $R$&  $G$& $B$\\
    \hline
         0&  0.0&  13&  8& 135\\
    \hline
         1&  0.2&  126&  3& 168\\
    \hline
         2&  0.4&  177&  42& 144\\
    \hline
         3&  0.6&  225&  100& 98\\
    \hline
         4&  0.8&  252&  166& 54\\
    \hline
         5&  1.0&  240&  249& 33\\
    \hline
    \end{tabular}
    \caption{Plasma anchor colors used to represent motivation in the simulator}
\end{table}

For $m\in[0,1]$, let $u=5m$, $k=\lfloor u\rfloor$ capped at 4, and $f=u-k$. I calculate the rendered color as:
\[
C(m)=(1-f)c_k+fc_{k+1}.
\]
I hold this mapping fixed across all agents, periods, scenarios, and simulation runs.

\subsection{Core Nodes, Commitment Arcs, and Contextual Public Nodes}

An outer ring identifies each modeled agent. The function $C(m_i(t))$ determines the node fill, while an arc with angular extent $2\pi L_i(t)$ represents commitment.

Inactive modeled agents are rendered with a desaturated or dashed style.

Public-body dots are contextual display elements outside the modeled node set $V$. They may be animated to illustrate reach but are excluded from all state updates and reported model statistics.

\subsection{Aggregate Readouts}
Mean motivation is
\[
\bar{m}(t)=\frac{1}{n}\sum_{i=1}^nm_i(t).
\]
For threshold $\tau$, motivational reach is
\[
\rho_\tau(t)=\frac{1}{n}|\{i:m_i(t)>\tau\}|.
\]
Cross-sectional motivational dispersion is
\[
s_m(t)=\sqrt{\frac{1}{n}\sum_{i=1}^n[m_i(t)-\bar{m}(t)]^2}
\]
I also report realized participation, expected active population, mean commitment, mean fatigue, motivational reach, terminal-window participation, and quorum-attainment or terminal-target-attainment rates.
\subsection{Ensemble Display}
Ensemble trajectories display the mean and, when available, empirical quantile bands across randomized realizations. For a scalar or indexed node quantity $Z_i(t)$,
\[
\bar{Z}_i(t)=\frac{1}{R}\sum_{r=1}^RZ_i^{(r)}(t).
\]
Node-level ensemble averages refer to indexed simulated agents or matched structural positions. 

In ensemble display mode, node states are shown as ensemble averages, while animated edges come from the realization whose terminal-window outcome is closest to the ensemble mean.

\section{Parameter, Initialization, and Reproducibility Tables}
Tables B1–B3 report the values used in the current implementation. “Fixed” means that the quantity remains constant across agents and realizations within the stated regime. “Randomized once” means that the quantity is sampled during initialization and then held fixed throughout that realization. “Dynamic” means that the model equations update the quantity over time.
For $X \sim \mathcal{N}(\mu, \sigma^2)$, I define the clipped-normal draw as:
\[
\text{CN}_{[a,b]}(\mu,\sigma)=\min\{b,\max\{a,X\}\}.
\]
The implementation draws one value from the stated normal distribution and then clips it to the stated interval. It does not redraw a value when the original draw falls outside that interval.

\begingroup
\small
\setlength{\tabcolsep}{3pt}
\renewcommand{\arraystretch}{1.08}
\begin{longtable}[c]{|c|c|c|c|c|c|}
    \caption*{Table B1. Global, Dynamic, and Control Parameters}
    \\
    \hline
    \textbf{Symbol} & \textbf{Meaning} & \makecell{\textbf{Allowed}\\\textbf{range}} & \makecell{\textbf{Collapse-prone}\\\textbf{setting}} & \makecell{\textbf{Favorable}\\\textbf{setting}} & \textbf{Status} \\ \hline
    \endhead
        \hline
         $n$&  \makecell{Number of\\modeled agents}&  $\mathbb{Z}^+$&  $36$&  $36$& Fixed\\
        \hline
         $T$&  \makecell{Number of\\simulated periods}&  $\mathbb{Z}^+$&  $18$&  $18$& Fixed\\
        \hline
         $N_{min}$&  \makecell{Minimum viable\\level of realized\\participation}&  $0, \ldots, n$& \makecell{$10$;\\sweep $2-12$}&  $4$& \makecell{Fixed\\within run}\\
        \hline
         $K$&  \makecell{Participation\\capacity}&  $>0$&  \makecell{$A: 14$;\\$C:20$}&  \makecell{$B: 18$}& \makecell{Fixed\\within run}\\
         \hline
         $\gamma_v$&  \makecell{Run-level\\visibility\\multiplier}&  $\ge1$&  \makecell{$A: 1.0; C: 3.0$; \\sweep $1.0$-$4.0$\\by $0.5$}& $B: 1.4$& \makecell{Fixed\\within run}\\
         \hline
         $\gamma_B(t)$&  \makecell{Time-dependent\\scaffolding\\multiplier}&  $\ge1$&  \makecell{$A: 1; C: 1.4$\\during window}&  $1$& \makecell{Dynamic\\input}\\
         \hline
         $s_B$&  \makecell{Below-quorum\\scaffolding support}&  $\ge0$&  $0.62$&  $0.62$& Fixed\\
         \hline
         $\lambda_R$&  \makecell{Event-pulse\\decay rate}&  $>0$&  $0.60$&  $0.60$& Fixed\\
         \hline
         $\kappa_R$&  \makecell{Event-input\\gain}&  $>0$&  $3.4$&  $3.4$& Fixed\\
         \hline
         $\{(\tau_k,A_k)\}$&  \makecell{Event\\schedule}&  $A_k\ge0$&  \makecell{Scenario\\specific}&  \makecell{Scenario\\specific}& Input\\
         \hline
         $\epsilon_0$&  \makecell{Safeguard against\\division by zero}&  $>0$&  $10^{-6}$&  $10^{-6}$& Fixed\\
         \hline
         $R$&  \makecell{Number of\\ensemble realizations}&  $\mathbb{Z}^+$&  \makecell{$250$ for trajectory\\ensembles; 20 per\\viability-visibility cell}&  \makecell{$250$ for trajectory\\ensembles; 20 per\\viability-visibility cell}& Fixed\\
         \hline
         $w$&  \makecell{Length of the\\terminal averaging\\window}&  $\mathbb{Z}^+$&  $4$&  $4$& Fixed\\
         \hline
         $A^*$&  \makecell{Terminal\\participation target}&  $[0, n]$&  $9$&  $9$& Fixed\\
         \hline
         $\Delta$&  \makecell{Credibility-history\\window}&  $\mathbb{Z}^+$&  $3$&  $3$& Fixed\\
         \hline
         $\pi_o,\pi_f$&  \makecell{Online/offline\\channel weights}&  \makecell{$[0, 1]$,\\sum to 1}&  $0.5, 0.5$&  $0.5, 0.5$& Fixed\\
         \hline
         $\mu$&  \makecell{Activity-dependent\\visibility update rate}&  $[0,1]$&  $0.28$&  $0.28$& Fixed\\
         \hline
         $\mu_{mix}$&  \makecell{Persistent-channel\\update rate}&  $[0,1]$&  $0.12$&  $0.12$& Fixed\\
         \hline
         $\eta_1, \eta_2$&  \makecell{Co-participation and\\broadcasting effects}&  $\ge0$&  $0.55, 0.45$&  $0.55, 0.45$& Fixed\\
         \hline
         $w_1^c, w_2^c$&  \makecell{Credibility\\weights}&  $\ge0$&  $0.5, 0.5$&  $0.5, 0.5$& Fixed\\
         \hline
         $\kappa_g$&  \makecell{Sender-signal\\saturation slope}&  $>0$&  2.0&  2.0& Fixed\\
         \hline
         $\eta_L$&  \makecell{Commitment\\protection\\of motivation}&  $\ge0$&  $0.93$&  $0.62$& Fixed\\
         \hline
         $\eta_F$&  \makecell{Fatigue effect on\\motivation decay}&  $\ge0$&  $0.55$&  $0.35$& Fixed\\
         \hline
         $\kappa_{SL}$&  \makecell{Saturation\\weakening\\of commitment}&  $[0,1]$&  $0.50$&  $0.50$& Fixed\\
         \hline
         $\lambda_L$&  \makecell{Commitment\\adjustment rate}&  $[0,1]$&  $0.37$&  $0.37$& Fixed\\
         \hline
         $\lambda_F$&  \makecell{Fatigue\\adjustment rate}&  $[0,1]$&  $0.30$&  $0.30$& Fixed\\
         \hline
         $\phi_1, \phi_2, \phi_3$&  \makecell{Commitment\\target weights}&  $\ge0$&  $0.68$, $0.34$, $0.20$&  $0.68$, $0.34$, $0.20$& Fixed\\
         \hline
         $\psi_1, \psi_2, \psi_3$&  \makecell{Fatigue\\target weights}&  $\ge0$&  $0.45$, $0.20$, $0.30$&  $0.45$, $0.20$, $0.30$& Fixed\\
         \hline
         $\gamma_1, \gamma_2, \gamma_3$&  \makecell{Dropout-hazard\\state coefficients}&  $\ge0$&  $1.6$, $1.6$, $1.3$&  $1.6$, $1.1$, $1.2$& Fixed\\
         \hline
         $\xi_1, \xi_2$&  \makecell{Reactivation\\coefficients}&  $\ge0$&  $1.4$, $1.5$&  $1.5$, $1.1$& Fixed\\
         \hline
         $\beta_E, \beta_F, \beta_R$&  \makecell{Affect\\couplings}&  \makecell{Stated model\\bounds}&  $0.15$, $0.30$, $0.30$&  $0.15$, $0.30$, $0.30$& Fixed\\
         \hline
         $\kappa_\Phi$&  \makecell{Stability-response\\coupling}&  $[0,1)$&  $0.3$&  $0.3$& Fixed\\
         \hline
         $\lambda_E, \lambda_S, \lambda_v$&  \makecell{Optional-state\\update rates}&  $[0,1]$&  $0.15$, $0.10$, $0.20$&  $0.15$, $0.10$, $0.20$& Fixed\\
         \hline
         $\kappa_{\text{var}}$&  \makecell{Variance-to-stability\\coefficient}&  $\ge0$&  18&  18& Fixed\\
         \hline
         $K_0$&  \makecell{Baseline capacity\\in cost function}&  $[0,n]$&  14&  14& Fixed\\
         \hline
         $\kappa_{E,0}, \kappa_{E,1}$&  \makecell{Event cost\\coefficients}&  $\ge0$&  $1.0$, $2.2$&  $1.0$, $2.2$& Fixed\\
         \hline
         $\kappa_v, \kappa_B, \kappa_K$&  \makecell{Visibility,\\scaffolding,\\and capacity costs}&  $\ge0$&  $0.9$, $1.4$, $0.0$&  $0.9$, $1.4$, $0.0$& Fixed\\
         \hline
         $\kappa_{EF}$&  \makecell{Affect adjustment\\to effective fatigue}&  $\ge0$&  $0.30$&  $0.30$& Fixed\\
         \hline
         $\zeta_E, \zeta_E'$&  \makecell{Positive and negative\\reinforcement effects\\on affect}&  $\ge0$&  $0.30$, $0.30$&  $0.30$, $0.30$& Fixed\\
         \hline
\end{longtable}
\endgroup
\begingroup
\small
\setlength{\tabcolsep}{3pt}
\renewcommand{\arraystretch}{1.08}
\begin{longtable}[c]{|c|c|c|c|c|c|}
    \caption*{Table B2a. Agent-Level Values: Core Initial States}
    \\
    \hline
    \textbf{Symbol} & \textbf{Meaning} & \textbf{Range} & \makecell{\textbf{Collapse-prone}\\\textbf{distribution}} & \makecell{\textbf{Favorable}\\\textbf{distribution}} & \textbf{Status} \\ \hline
    \endhead
         \hline
         $m_i(0)$&  Initial motivation&  $[0,1]$&  $\text{CN}_{[0,1]}(0.28,0.10)$&  $\text{CN}_{[0,1]}(0.33,0.10)$& \makecell{Randomized\\once}\\
         \hline
         $L_i(0)$&  \makecell{Initial\\commitment}&  $[0,1]$&  $\text{CN}_{[0,1]}(0.10,0.05)$&  $\text{CN}_{[0,1]}(0.12,0.06)$& \makecell{Randomized\\once}\\
         \hline
         $F_i(0)$&  \makecell{Initial\\fatigue}&  $[0,1]$&  $\text{CN}_{[0,1]}(0.35,0.10)$&  $\text{CN}_{[0,1]}(0.28,0.10)$& \makecell{Randomized once}\\
         \hline
         $\bar{E}_i$&  \makecell{Affect\\baseline}&  $[0,1]$&  $\text{CN}_{[0,1]}(0.45,0.12)$&  $\text{CN}_{[0,1]}(0.22,0.07)$& \makecell{Randomized\\once}\\
         \hline
         $E_i(0)$&  \makecell{Initial\\affect}&  $[0,1]$&  $E_i(0)=\bar{E}_i$&  $E_i(0)=\bar{E}_i$& \makecell{Initialized as\\$E_i(0)=\bar{E}_i$}\\
         \hline
         $S_i(0)$&  \makecell{Initial\\saturation}&  $[0,1]$&  $0$&  $0$& \makecell{Fixed}\\
         \hline
         $\Phi_i(0)$&  \makecell{Initial\\stability gain}&  $[0,1]$&  $1$&  $1$& \makecell{Fixed}\\
         \hline
         $act_i(0)$&  \makecell{Initial activity\\status}&  \makecell{\{0,1\} or\\$[0,1]$}&  $1$&  $1$& \makecell{Fixed}\\
         \hline
         $\bar{m}_i(0)$&  \makecell{Initial moving\\mean}&  [0,1]&  $\bar{m}_i(0)=m_i(0)$&  $\bar{m}_i(0)=m_i(0)$& \makecell{Initialized}\\
         \hline
         $\mathrm{var}_i(0)$&  \makecell{Initial moving\\variance}&  $[0,1]$&  $0$&  $0$& \makecell{Fixed}\\
         \hline
\end{longtable}
\endgroup
\begingroup
\small
\setlength{\tabcolsep}{3pt}
\renewcommand{\arraystretch}{1.08}
\begin{longtable}[c]{|c|c|c|c|c|c|}
    \caption*{Table B2b. Agent-Level Values: Behavioral Parameters}
    \\
    \hline
    \textbf{Symbol} & \textbf{Meaning} & \textbf{Range} & \makecell{\textbf{Collapse-prone}\\\textbf{distribution}} & \makecell{\textbf{Favorable}\\\textbf{distribution}} & \textbf{Status} \\ \hline
    \endhead
         \hline
         $\theta_i^a$&  \makecell{Participation\\threshold}&  $[0,1]$&  $\text{CN}_{[0.1,0.95]}(0.60,0.08)$&  $\text{CN}_{[0.1,0.95]}(0.56,0.10)$& \makecell{Randomized\\once}\\
         \hline
         $\theta_i^p$&  \makecell{Broadcasting\\threshold}&  $[0,1]$&  $\text{CN}_{[0.1,0.98]}(0.70,0.08)$&  $\text{CN}_{[0.1,0.98]}(0.70,0.12)$& \makecell{Randomized\\once}\\
         \hline
         $\kappa_i^a$&  \makecell{Participation-\\response slope}&  $>0$&  9&  8& \makecell{Fixed}\\
         \hline
         $\kappa_i^p$&  \makecell{Broadcasting-\\response slope}&  $>0$&  11&  9& \makecell{Fixed}\\
         \hline
         $\alpha_{0,i}$&  \makecell{Baseline\\motivation\\decay}&  $[0,1]$&  $\text{CN}_{[0.02,0.95]}(0.45,0.10)$&  $\text{CN}_{[0.02,0.95]}(0.28,0.08)$& \makecell{Randomized\\once}\\
         \hline
         $w_i^P$&  \makecell{Peer-input\\susceptibility}&  $\ge0$&  $\text{CN}_{[0,1.2]}(0.28,0.10)$&  $\text{CN}_{[0,1.2]}(0.40,0.15)$& \makecell{Randomized\\once}\\
         \hline
         $b_i$&  \makecell{Sender\\spreading\\propensity}&  $\ge0$&  $\text{CN}_{[0,1]}(0.30,0.10)$&  $\text{CN}_{[0,1]}(0.46,0.10)$& \makecell{Randomized\\once}\\
         \hline
         $\rho_i$&  \makecell{Peer-input\\saturation\\coefficient}&  $\ge0$&  $\text{CN}_{[0.3,3]}(1.30,0.20)$&  $\text{CN}_{[0.3,3]}(1.50,0.25)$& \makecell{Randomized\\once}\\
         \hline
         $\beta_i$&  \makecell{Maximum\\productive\\reinforcement}&  $\ge0$&  $\text{CN}_{[0,3]}(0.26,0.08)$&  $\text{CN}_{[0,3]}(0.34,0.11)$& \makecell{Randomized\\once}\\
         \hline
         $\delta_i$&  \makecell{Reinforcement\\saturation rate}&  $>0$&  $2.0$&  1.25& \makecell{Fixed}\\
         \hline
         $\chi_i$&  \makecell{Nonparticipation\\penalty}&  $\ge0$&  $\text{CN}_{[0,0.5]}(0.06,0.03)$&  $\text{CN}_{[0,0.5]}(0.05,0.03)$& \makecell{Randomized\\once}\\
         \hline
         $\chi_i^{\text{viab}}$&  \makecell{Below-quorum\\penalty}&  $\ge0$&  $0.016$&  $0.040$& \makecell{Fixed}\\
         \hline
         $h_{0,i}$&  \makecell{Baseline\\dropout hazard}&  $[0,0.3]$&  $\text{CN}_{[0,0.3]}(0.06,0.02)$&  $\text{CN}_{[0,0.3]}(0.05,0.02)$& \makecell{Randomized\\once}\\
         \hline
         $V_i^{\text{ev}}$&  \makecell{Event\\visibility}&  $[0,1]$&  $0.70$&  $0.72$& \makecell{Fixed}\\
         \hline
         $V_{i,0}^{\text{ex}}$&  \makecell{Baseline\\expressive reach}&  $[0,1]$&  $\text{CN}_{[0,1]}(0.25,0.10)$&  $\text{CN}_{[0,1]}(0.60,0.13)$& \makecell{Randomized\\once}\\
         \hline
         $\zeta_i^0$&  \makecell{Baseline\\trust}&  [0,1]&  $0.3$&  $0.3$& \makecell{Fixed}\\
         \hline
         $\zeta_i^1$&  \makecell{Similarity\\weight}&  $\ge0$&  $0.40$&  $0.40$& \makecell{Fixed}\\
         \hline
         $\zeta_i^2$&  \makecell{Credibility\\weight}&  $\ge0$&  $0.30$&  $0.30$& \makecell{Fixed}\\
         \hline
         $s_0$&  \makecell{Global saturation\\half-load parameter}&  $>0$&  $2.0$&  $2.0$& \makecell{Fixed}\\
         \hline
\end{longtable}
\endgroup
\begingroup
\small
\setlength{\tabcolsep}{3pt}
\renewcommand{\arraystretch}{1.08}
\begin{longtable}[c]{|c|c|}
    \caption*{Table B2c. Agent-Level Values: Inputs and Shocks}
    \\
    \hline
    \textbf{Symbol} & \textbf{Exact Implementation} \\ \hline
    \endhead
         \hline
         $c_i(t)$&  Drawn once as $\text{CN}_{[0.1,1]}(0.55,0.12)$, then constant over time\\
         \hline
         $T_i(t)$&  $1$ for every agent and period\\
         \hline
         $\bar{y}_i$&  $1$\\
         \hline
         $s_i^{\text{role}}(t)$&  $1$\\
         \hline
         $r^+(t)$&  $0.10$\\
         \hline
         $\epsilon_i(t)$&  \makecell{$0$ in deterministic mode;\\$\text{Uniform}(-0.01,0.01)$ in stochastic mode}\\
         \hline
         $\kappa_E\xi_i^E(t)$&  \makecell{$0$ in deterministic mode;\\$\text{Uniform}(-0.05,0.05)$ in stochastic mode}\\
         \hline
\end{longtable}
\endgroup
\begingroup
\small
\setlength{\tabcolsep}{3pt}
\renewcommand{\arraystretch}{1.08}
\begin{longtable}{|c|c|}
    \caption*{Table B3. Network Generation, Initialization, and Reproducibility Settings}
    \\
    \hline
    \textbf{Component} & \textbf{Exact Implementation} \\ \hline
    \endhead
         \hline
         Number of communities $C$&  $C=\max\{2,\text{round}(n/g)\};$ reference $C=5$\\
         \hline
         Community-size rule&  \makecell{For $n=36$, $g=8$,\\the community sizes are $8$, $7$, $7$, $7$, $7$}\\
         \hline
         Community assignment&  $\text{comm}_i=i\mod{C}$\\
         \hline
         Angular position&  \makecell{Community-sector center with an added draw\\from $\text{Uniform}(-0.8\pi/C,0.8\pi/C)$}\\
         \hline
         \makecell{Within-community\\tie rule}&  \makecell{Each receiver selects $4+U$ most similar senders\\from the same community. $U \sim \text{DiscreteUniform}\{0,1,2\}$}\\
         \hline
         \makecell{Cross-community\\bridge rule}&  \makecell{Each receiver selects one or two uniformly sampled senders\\from other communities, with probabilities $0.5,0.5$}\\
         \hline
         Edge direction&  A selected sender $j$ forms directed edge $j \rightarrow i$\\
         \hline
         Self-loops&  Prohibited\\
         \hline
         Reciprocal ties&  Allowed but generated independently\\
         \hline
         $\text{sim}_{ij}$& $\exp(-d_{ij}^2/0.6)$\\
         \hline
         Duplicate directed edges&  Prohibited\\
         \hline
         $\text{online}_{ij}(0)$&  Equal to $v_{ij}^{\text{mix}}(0)$\\
         \hline
         $\text{offline}_{ij}(0)$&  Equal to $v_{ij}^{\text{mix}}(0)$\\
         \hline
         $\pi_o,\pi_f$& $0.5, 0.5$\\
         \hline
         $v_{ij}(0)$&  $\clip[0.5v_{ij}^{\text{mix}}(0)+0.1]$\\
         \hline
         $\bar{a}_j^{\text{pre}}$&  $0$\\
         \hline
         $\text{cred}_j(0)$&  $0.5L_j(0)$\\
         \hline
         $s_{ij}(0)$& $\clip[0.30+0.40\text{sim}_{ij}+0.30\text{cred}_j(0)]$\\
         \hline
         $\omega_{ij}(0)$&  $\clip[v_{ij}(0)s_{ij}(0)]$\\
         \hline
         \makecell{Additional\\acquaintance ties}&  $0$ in all reported reference experiments\\
         \hline
         \makecell{Network structure}&  Resampled for each realization\\
         \hline
         \makecell{Heterogeneous\\parameters}&  Resampled for each realization\\
         \hline
         \makecell{Initial latent states}&  Resampled for each realization\\
         \hline
         Master seed&  11\\
         \hline
         \makecell{Ensemble seed sequence}&  $q_r=11+9973r$, $r=0,\ldots, R-1$\\
         \hline
         Language&  Vanilla JavaScript, ES2021-compatible\\
         \hline
         Interface technology&  HTML5 Canvas\\
         \hline
         External packages&  None\\
         \hline
         \makecell{Current SHA-256 hash\\of the implementation file}&  601d98427e75a058a499f0e269713e10a9c543ceda25c1bc4f0929746dfa6fb5\\
         \hline
\end{longtable}
\endgroup

\subsubsection*{Disclosure of Interests.}
The author has no competing interests to declare that are relevant to the content of this article.

\section*{References}

\noindent\hangindent=0.5in \hangafter=1
Grimm, V., Railsback, S. F., Vincenot, C. E., et al. (2020). The ODD protocol for describing agent-based and other simulation models: A second update to improve clarity, replication, and structural realism. Journal of Artificial Societies and Social Simulation, 23(2), 7. \\
https://doi.org/10.18564/jasss.4259

\end{document}
